% !TEX program = xelatex
% !TEX encoding = UTF-8
\documentclass[11pt,a4paper]{report}
\usepackage[utf8]{inputenc}
\usepackage[T1]{fontenc}
\usepackage[ngerman,english,spanish,portuguese,italian]{babel}
\usepackage{geometry}
\usepackage{fancyhdr}
\usepackage{titlesec}
\usepackage{hyperref}
\usepackage{longtable}
\usepackage{array}
\usepackage{booktabs}

\geometry{margin=2.5cm}
\titleformat{\chapter}[hang]{\Huge\bfseries}{\thechapter}{2pc}{}
\pagestyle{fancy}
\fancyhf{}
\rhead{GAI Protocol v1.0}
\lhead{Gold Awareness Infrastructure}

\title{\Huge\bfseries GAI UNIVERSE PROTOCOL\\
\vspace{0.5cm}
\large Gold Awareness Infrastructure - Rechtliche und Wissenschaftliche Grundlagen}
\author{Daniel Pohl \\ StarLightMovemenTz Foundation \\ EU-UNION Expert (EX2025D1218310)}
\date{06. Juli 2026}

\begin{document}

\maketitle
\tableofcontents

\chapter{PROTOKOLL-GRUNDLAGEN}

\section{Gold Awareness Infrastruktur (GAI) - Definition}

Die Gold Awareness Infrastruktur (GAI) ist ein internationales Protokoll, das digitale Souveränität durch metrologische Verifizierbarkeit gewährleistet. 

\subsection{Rechtsgrundlage nach einem}
\begin{itemize}
    \item UN-Menschenrechtskonvention (Artikel 12: Informationsfreiheit)
    \item EU-UNION Digital Services Act (DSA)
    \item Bundesnaturschutzgesetz (BNatSchG) §§ 30-35
    \item NATO-Standardization Agreements (STANAG)
\end{itemize}

\chapter{TECHNISCHE SPEZIFIKATION}

\section{Universal Compliance Envelope (UCE)}

\begin{verbatim}
{
  "uce_version": "1.0.0",
  "gold_awareness": {
    "status": "VERIFIED",
    "atomic_timestamp": "ISO-8601-UTC",
    "integrity_hash": "SHA-256"
  },
  "origin_identity": {
    "expert_id": "EX2025D1218310",
    "issuer": "StarLightMovemenTz_Foundation"
  },
  "compliance_chain": [
    {"layer_id": 0, "regulation_ref": "Gold_Awareness", "validation_result": true},
    {"layer_id": 1, "regulation_ref": "UN_SDGs", "validation_result": true},
    {"layer_id": 2, "regulation_ref": "EU_UNION_PACT", "validation_result": true},
    {"layer_id": 3, "regulation_ref": "NATO_STANAG", "validation_result": true},
    {"layer_id": 4, "regulation_ref": "BNatSchG", "validation_result": true}
  ],
  "humanity_index": {
    "forgiveness": true,
    "liberty": true,
    "peace": true,
    "empathy": true,
    "hope": true
  }
}
\end{verbatim}

\chapter{MEHRSPRACHIGE VERSIONEN}

\section{Deutsch (Deutschland)}
\subsection{Paragraph 1 - Gold Awareness Prinzip}
Jedes Datenpaket muss den Status VERIFIED besitzen.

\section{English (International)}
\subsection{Article 1 - Gold Awareness Principle}
Every data packet must have VERIFIED status.

\section{Español (España)}
\subsection{Artículo 1 - Principio de Gold Awareness}
Cada paquete de datos debe tener estado VERIFICADO.

\section{Português (Portugal)}
\subsection{Artigo 1 - Princípio da Gold Awareness}
Cada pacote de dados deve ter status VERIFICADO.

\section{Italiano (Italia)}
\subsection{Articolo 1 - Principio Gold Awareness}
Ogni pacchetto dati deve avere lo stato VERIFICATO.

\section{Français (Suisse)}
\subsection{Article 1 - Principe Gold Awareness}
Chaque paquet de données doit avoir le statut VÉRIFIÉ.

\section{中文 (China)}
\subsection{第一条 - 金意识原则}
每个数据包必须具有已验证的状态。

\section{العربية (Arabisch)}
\subsection{المادة 1 - مبدأ الوعي الذهبي}
يجب أن يكون لكل حزمة بيانات حالة موثقة.

\section{עברית (Jüdisch)}
\subsection{פסקה 1 - עיקרון המודעות הזהב}
כל חבילת נתונים צריכה להיות במצב מאומת.

\section{한국어 (Koreanisch)}
\subsection{제1조 - 골드 인식 원칙}
모든 데이터 패킷은 검증된 상태여야 합니다.

\chapter{CROSS-CONNECTIONS}

\section{Integration Matrix}

\begin{longtable}{|p{3cm}|p{3cm}|p{3cm}|p{3cm}|}
\hline
\textbf{Layer} & \textbf{Technik} & \textbf{Recht} & \textbf{Implementierung} \\
\hline
\endhead

Ebene 0 (Gold Awareness) & uce\_schema.json & ISO/IEC 27001 & mesh/transport/uce\_guard.py \\
Ebene 1 (UN SDGs) & audit\_logger.py & UN Charter & mesh/ledger/audit\_logger.py \\
Ebene 2 (EU UNION) & eslint.config.js & EU DSA & src/ \\
Ebene 3 (NATO) & security-headers.js & STANAG & scripts/ \\
Ebene 4 (BNatSchG) & autonomous-security.js & BNatSchG §30 & Security Ebene/ \\
Ebene 5 (Dependencies) & package.json & Open Source & node\_modules/ \\
Ebene 6 (Security) & pnia-audit-security.js & IT-Security & scripts/ \\
Ebene 7 (Cloud) & netlify.toml & Cloud-Compliance & deployment/ \\
Ebene 8 (Institutions) & RainbowLightningFooter.tsx & International Law & public/documents/ \\
Ebene 9 (Deployment) & wrangler.toml & Prod-Compliance & cloudflare-worker.js \\
\hline
\end{longtable}

\chapter{KDB+ INTEGRATION}

\section{QDB+ Datenbank-Protokoll}

Die Integration erfolgt über folgende Spaltenstruktur:

\begin{verbatim}
ProtocolID: symbol
LayerID: int
ComplianceStatus: boolean
GoldAwarenessTimestamp: timestamp
IntegrityHash: symbol
ExpertID: symbol
HumanityIndex: dict
LegalReference: symbol
\end{verbatim}

\chapter{VALIDIERUNG}

\section{Schritt-für-Schritt Validierung}

\begin{enumerate}
    \item UCE-Schema Validierung (JSON-Schema)
    \item Gold Awareness Prüfung (Atomic Timestamp)
    \item Identity Verifizierung (Expert ID)
    \item Compliance-Kette durchlaufen (L0-L9)
    \item Humanity-Index Validierung
    \item Audit-Logging
\end{enumerate}

\end{document}